% Generated draft: requires mathematical and source review.
\documentclass[11pt]{article}
\usepackage[margin=25mm]{geometry}
\usepackage{fontspec}
\setmainfont{Latin Modern Roman}
\usepackage{amsmath,amssymb,mathtools,unicode-math}
\setmathfont{Latin Modern Math}
\usepackage{enumitem,needspace,xurl,hyperref}
\setlength{\parindent}{0pt}
\setlength{\parskip}{5pt}
\setlength{\emergencystretch}{3em}
\newcommand{\N}{\mathbb{N}}
\newcommand{\Z}{\mathbb{Z}}
\newcommand{\Q}{\mathbb{Q}}
\newcommand{\R}{\mathbb{R}}
\newcommand{\C}{\mathbb{C}}
\newcommand{\F}{\mathbb{F}}
\newcommand{\A}{\mathbb{A}}
\newcommand{\PP}{\mathbb{P}}
\newcommand{\E}{\mathbb{E}}
\newcommand{\eps}{\varepsilon}
\newcommand{\dd}{\,\mathrm{d}}
\begin{document}
\section*{UnsolvedMath ID: 30000596}
\textbf{Source problem code:} catalogue entry
\subsection{Definitions and mathematical statement}
Let $(M,g)$ be a simple compact Riemannian $n$-manifold with boundary. In the supplied definition, simplicity means strictly convex boundary, a unique geodesic segment joining every pair of points, and no conjugate points. For a Riemannian manifold with boundary, $d_M$ is its intrinsic path distance and $\operatorname{vol}_n(M)$ its Riemannian volume.

In the usual filling category, let $(M',g')$ be any compact orientable Riemannian $n$-manifold with a fixed boundary identification $\iota:\partial M'\to\partial M$. The conjecture asserts
\[
\left[\forall x,y\in\partial M',\quad
d_{M'}(x,y)\geq d_M(\iota(x),\iota(y))\right]
\quad\Longrightarrow\quad
\operatorname{vol}_n(M')\geq\operatorname{vol}_n(M).
\]
Furthermore, under the same boundary-distance hypothesis, equality of volumes should imply existence of an isometry
\[
F:(M',g')\longrightarrow(M,g),\qquad F|_{\partial M'}=\iota.
\]
Thus the equality case respects the specified boundary identification. The competitor need not be simple or diffeomorphic to $M$, and the hypothesis compares intrinsic distances of boundary points through the fillings, not just the metrics intrinsic to their boundaries.

The extract explicitly imposes orientability on $M'$ but does not separately spell out its compactness and dimension; the displayed formulation uses the standard compact, same-dimensional Riemannian filling interpretation. Any broader class of competitors would need its own volume and filling conventions.
\subsection{Short English statement}
A simple compact Riemannian manifold should minimize volume among orientable fillings whose distances between corresponding boundary points are no smaller. Equality should force an isometry fixing the prescribed boundary identification.
\subsection{Sources}
\begin{itemize}
\item Supplied report: \emph{Geometrie} (2007), Oberwolfach Reports, DOI 10.4171/owr/2006/45, pp. 2699--2756; proposed year 2006. \url{https://doi.org/10.4171/owr/2006/45}.
\item Supplied reference: Ivanov, \emph{Volume comparison via boundary distances}, \url{https://arxiv.org/abs/1004.2505}.
\item Local UnsolvedMath registry, permanent ID 30000596, catalogue entry; created 25 August 2026, updated 16 September 2026. Its corrected\_verified note says the equality case and definition of simplicity were retained. The supplied 21 August 2026 literature triage labels the general conjecture open, reporting metrics close to flat and other special classes rather than all simple manifolds. This dated assessment has not been independently checked here.
\end{itemize}
\end{document}
